% English translation of the verified Chinese baseline. Structural identifiers must remain synchronized.
\section*{Introduction}
\label{sec:introduction}

Magnetically actuated flexible catheters in a microgravity environment are simultaneously subjected to magnetic actuation, background-flow effects, vessel-wall contact, and boundary manipulation. Their dynamics therefore combine large rotations, strong geometric nonlinearity, bending--torsion coupling, and open-system ports. If the continuous model, finite-dimensional reduction, data representation, and long-time stability analysis use inconsistent variables or energy conventions, model error, learning error, and numerical discretization error cannot be compared within a common theoretical framework. Consistent mathematical interfaces are therefore required across continuum mechanics, structure-preserving discretization, and verifiable machine learning.

Geometrically exact beam and rod theories provide a unified kinematic and variational foundation for flexible structures undergoing three-dimensional large deformation, while magnetic, fluid, and contact models describe the surrounding multiphysics environment\cite{Simo1985FiniteStrainBeam,Antman2005Elasticity,Morison1950,COMSOLBeamInterface2025,COMSOLGeometricNonlinearity2025,COMSOLSectionStiffness2025,COMSOLWeakFormPDE2025,COMSOLWeakContribution2025,Sola2018MicroLie,Jackson1998Electrodynamics,Abbott2020MagneticRobotics,Newman1977MarineHydrodynamics,Faltinsen1990SeaLoads,Johnson1985ContactMechanics,HuntCrossley1975Restitution,BrownMcPhee2016Friction,MarsdenRatiu1999MechanicsSymmetry,HairerLubichWanner2006GNI}. However, retaining translation and cross-sectional orientation jointly, enforcing mutually exclusive energy accounting for conservative and nonconservative loads, and obtaining a finite-dimensional canonical system suitable for Hamiltonian learning and symplectic analysis require a consistent construction of the weak form, Galerkin discretization, constraint reduction, and Legendre transformation.

Hamiltonian neural networks introduce a conservative inductive bias through the symplectic gradient of a real-valued energy function. Dissipative and port-based networks extend structured identification to forced and dissipative systems, whereas complex-valued networks provide structured representations for amplitude--phase coupling\cite{Greydanus2019HNN,Zhong2020Dissipative,Desai2021PHNN,Chen2022KAMSLT,Trabelsi2018DeepComplex,Hirose2012CVNN}. A complex-valued internal representation, however, does not imply that the physical Hamiltonian may be an arbitrary complex-valued function, nor can dissipation and external forcing be generated unconditionally by a single complex scalar. Moreover, a multi-branch architecture does not guarantee identifiability by construction; it must be validated using a known physical backbone, port-power supervision, paired operating conditions, and ablation studies.

KAM theory and symplectic geometric algorithms provide rigorous tools for analyzing quasiperiodic structures in nearly integrable Hamiltonian systems and their numerical maps\cite{FengQin2010,Shang1999KAMAlgorithms,Shang2000KAMMappings,Haro2016Parameterization,Arnold1989ClassicalMechanics,Coquinot2020Metriplectic,Moser1962InvariantCurves,Kuksin1993InfiniteDimensionalKAM,Poschel2001KAMLecture}. Classical continuous KAM theory, KAM theory for symplectic maps, two-dimensional twist-map theory, and infinite-dimensional KAM theory concern different mathematical objects. Neither a training loss nor a trajectory residual can directly replace the required smallness condition in a high-order or analytic function norm. Learning errors must therefore be measured together with physical perturbations, coordinate-transformation errors, Galerkin truncation, and time-discretization errors in compatible norms.

The main contributions of this work are summarized as follows.

(1) A multiphysics geometrically exact catheter model is established that retains centerline translation, cross-sectional orientation, independent torsion, and rotary inertia, while magnetic, fluid, contact, and boundary ports are organized through a unified energy--power identity.

(2) A joint translational--rotational Galerkin discretization, local reduction of the Kirchhoff constraints, and a consistent Legendre transformation are used to construct a finite-dimensional canonical Hamiltonian system compatible with COMSOL data interfaces.

(3) A structured neural network with complex-valued internal representations and a real-valued Hamiltonian readout is developed, with separate branches for conservative corrections, passive dissipation, external inputs, and general residuals.

(4) Interfaces are established among statistical learning error, high-order regularity certificates, symplectic-discretization error, and KAM conditions, yielding separate analysis routes for continuous flows, symplectic maps, and a posteriori invariant-torus validation.

The numerical study comprises the low-dimensional P0 architecture comparison, the three-trajectory P1/P1.5/P1.6 structural ablations, the COMSOL--Galerkin bridge, and diagnosis of the canonical data contract. The present three-level time-step analysis and \(1\,\mathrm{s}\) extended case show that this contract meets the data-generation requirements for subsequent multi-amplitude sentinel experiments. The planned 96-trajectory data set, full-scale P2 training, open-port identification, high-fidelity model comparison, and numerical KAM analysis have not yet been completed; consequently, the present pilot results are not used to claim full-scale model superiority, invariant tori, or an engineering safe domain.

The remainder of the paper is organized as follows. \Cref{sec:rod-model,sec:galerkin-discretization} develop the continuous multiphysics model and the joint Galerkin discretization. \Cref{sec:hamiltonian-structure} constructs the finite-dimensional Hamiltonian structure and canonical interfaces. \Cref{sec:chnn-model} presents the complex-representation structured HNN and its statistical error analysis. \Cref{sec:symplectic-kam} discusses symplectic discretization, Birkhoff structures, and invariant-torus validation. Numerical results are reported in \Cref{sec:numerical-results}, applicability and limitations are discussed in \Cref{sec:discussion}, and \Cref{sec:conclusion} concludes the paper.
